\documentclass[11pt,a4paper]{article}
\usepackage[margin=2.5cm]{geometry}
\usepackage[T1]{fontenc}
\usepackage[utf8]{inputenc}
\usepackage{lmodern}
\usepackage{microtype}
\usepackage[hidelinks]{hyperref}
\usepackage{parskip}

\title{KherveTeX (kherveDOC): A WYSIWYG LaTeX Document Editor\\with Git-Based Version History}
\author{Gwilherm Kerhervé\thanks{ORCID: 0000-0002-6449-1828, \texttt{g.kerherve@imperial.ac.uk}}\\
\small Department of Materials, Imperial College London, London SW7 2AZ, United Kingdom}
\date{10 July 2026 --- Version 0.153}

\begin{document}
\maketitle

\begin{abstract}
KherveTeX (repository name \texttt{kherveDOC}) is an open-source WYSIWYG document editor that produces publication-quality PDFs via LaTeX. Documents are edited on a Word-style paginated surface backed by a typed document model; on every edit the model is serialised to LaTeX and compiled with the \texttt{tectonic} engine, giving a live PDF preview. The editor supports inline and display mathematics, figures, tables, citations, cross-references, footnotes, and import of \texttt{.docx} and \texttt{.tex} files. Every document carries its own Git repository with auto-commit on save and optional push to a remote. KherveTeX is written in Python with PySide6 and released under the GNU General Public License v3.0 at \url{https://github.com/gkerherve/kherveDOC}.
\end{abstract}

\section{Statement of need}

Scientific writing is split between two workflows. Word processors are approachable but offer weak support for mathematics, cross-referencing and consistent typography; LaTeX produces excellent documents but presents a steep learning curve that excludes many students and experimentalists. Bridging tools exist --- LyX~\cite{lyx} and TeXmacs~\cite{texmacs} pioneered structured WYSIWYG editing, and Overleaf moved LaTeX into the browser --- yet there remains room for a lightweight, native editor that (i) looks and behaves like a familiar word processor, (ii) keeps human-readable LaTeX as its archival output, and (iii) records the full history of a document automatically.

KherveTeX addresses this gap. It was developed to let researchers and students who have never written LaTeX produce properly typeset reports, theses and manuscripts, while still emitting clean \texttt{.tex} sources that collaborators or journals can consume directly. The built-in Git history removes the ``\texttt{report\_final\_v3\_REALLY\_final.docx}'' problem: every save is a commit, and the history viewer allows any previous state to be inspected or restored.

\section{Functionality}

The editing surface is a Word-style page card rendered at real A4, Letter or Legal dimensions, with dashed page-break lines at every layout pagination. Users write with a paragraph-style picker (Body, Title, Headings 1--5) and inline marks (bold, italic, underline, strikethrough, code, small caps, sub/superscript). Structured inserts include inline and display mathematics (typed as LaTeX), bulleted and numbered lists, hyperlinks, footnotes, citations (\texttt{cite}/\texttt{citep}/\texttt{citet}), cross-references (\texttt{ref}/\texttt{eqref}/\texttt{pageref}), captioned and labelled figures, tables, page breaks, horizontal rules, and a raw-LaTeX escape hatch for anything not covered by the model.

Three synchronised tabs show the same document as (1) the formatted WYSIWYG view, (2) live, read-only, syntax-highlighted LaTeX source, and (3) the compiled PDF preview. The native format is \texttt{.ktex.json}, a JSON serialisation of the document model, saved alongside a sibling \texttt{.tex}; a bundled \texttt{.ktexz} ZIP format embeds images. Import is supported from \texttt{.tex} (a regex-based subset parser with raw-LaTeX fallback) and from \texttt{.docx} via python-docx, including extraction of embedded images. Export targets \texttt{.tex} and \texttt{.pdf}.

\section{Implementation}

KherveTeX follows a strict model-driven architecture: a typed document tree (Python dataclasses) is the single source of truth. The editor reads and writes the model; a serialiser converts the model to LaTeX; the \texttt{tectonic} engine~\cite{tectonic} compiles it; and PyMuPDF~\cite{pymupdf} rasterises the resulting pages into the preview pane. No component edits LaTeX strings directly. Version control uses pygit2~\cite{pygit2}: each document folder is initialised as its own Git repository, saving auto-commits the human-readable \texttt{.tex}, and pushes to \texttt{origin} on a best-effort basis. The GUI is built with PySide6 (Qt~6). The test suite covers the document model, the LaTeX serialiser and the importers.

\section{Availability}

Source code, a Windows installer and documentation are available at \url{https://github.com/gkerherve/kherveDOC} under the GPL-3.0 license. KherveTeX runs on Python~3.10+ with PySide6, PyMuPDF, pygit2 and the tectonic binary. It is part of the Kherve family of open-source scientific tools, which includes KherveFitting~\cite{khervefitting}, KherveBook and KhervePDF.

\begin{thebibliography}{9}
\bibitem{lyx} The LyX Team, \emph{LyX --- The Document Processor}, \url{https://www.lyx.org}.
\bibitem{texmacs} J.~van der Hoeven, \emph{GNU TeXmacs: A free, structured, wysiwyg and technical text editor}, Cahiers GUTenberg 39--40 (2001) 39--50.
\bibitem{tectonic} P.~K.~G. Williams et al., \emph{Tectonic: a modernized, complete, self-contained TeX/LaTeX engine}, \url{https://tectonic-typesetting.github.io}.
\bibitem{pymupdf} J.~X. McKie, \emph{PyMuPDF: Python bindings for MuPDF}, \url{https://github.com/pymupdf/PyMuPDF}.
\bibitem{pygit2} \emph{pygit2: Python bindings for libgit2}, \url{https://www.pygit2.org}.
\bibitem{khervefitting} G.~Kerhervé, D.~J. Payne, \emph{KherveFitting: An Open Source Software for Fitting X-Ray Photoelectron Spectroscopy Data}, Surf. Interface Anal. (2026), doi:10.1002/sia.70032.
\end{thebibliography}

\end{document}
